% Desarrollo reunido de DESARROLLO_MATEMATICO, apartados 38--40 y 50.
\subsubsection{Índice de capacidad, vacancias e inducción del régimen trítico}
\label{vac-capacidad-trit}

La sucesión de vacancias está vinculada al refinamiento mediante una
identidad de incrementos. Para mostrarla se fija el origen de fase
\(\theta=0\). Sean
\[
 \rho=\log_{10}9=1-\delta,\qquad
 \mathcal C_t=\lfloor(t+1)\rho\rfloor\quad(t\geq0).
\]
El índice \(\mathcal C_t\) expresa la capacidad decimal del refinamiento
en la normalización que compara un bloque ternario de \(6\) posiciones,
con \(729\) posibilidades, y un bloque decimal de \(3\) posiciones, con
\(1000\) posibilidades. En efecto,
\(\log_{1000}729=\log_{10}9=\rho\). Se reserva la letra \(K\) para
el registro dodecafásico que se construirá posteriormente.

\begin{proposition}[Complementariedad de capacidad y vacancia]
\label{prop:capacidad-vacancia}
Para \(t\geq1\),
\begin{equation}
 \mathcal C_t-\mathcal C_{t-1}=1-z_t(0).
 \label{eq:capacidad-vacancia}
\end{equation}
En consecuencia, si \(0\leq a<b\),
\[
 \mathcal C_b-\mathcal C_a+
 \sum_{t=a+1}^{b}z_t(0)=b-a.
\]
\end{proposition}
\begin{proof}
La irracionalidad de \(\delta\) da
\[
 \mathcal C_t
 =t+1-\lceil(t+1)\delta\rceil
 =t-\lfloor(t+1)\delta\rfloor.
\]
La resta de dos términos consecutivos produce
\eqref{eq:capacidad-vacancia}. La suma sobre el intervalo es telescópica.
\end{proof}

Cada transición aumenta en una unidad el contador de bloques. Si
\(z_t=0\), aumenta también la capacidad; si \(z_t=1\), la capacidad se
mantiene mientras se incorpora una transición a la historia. La vacancia
es aquí un evento de retención de capacidad, con una posición y un
antecedente determinados. El balance anterior conserva exactamente la
longitud entre capacidad adquirida y vacancias registradas. La fase
observada \(g_t^{\rm cap}=[\mathcal C_t]_9\) satisface
\[
 g_t^{\rm cap}-g_{t-1}^{\rm cap}\equiv1-z_t\pmod9.
\]
Esta fase de capacidad y la fase cronológica \([t]_9\) son dos
observaciones diferentes. Una retención local de capacidad y una vuelta
completa de la holonomía tienen mecanismos distintos, aunque ambas
permitan conservar la observación de fase con modificación del estado.

La relación con TRIT puede seguirse en las propias posiciones de vacancia.
Para \(j\geq1\), defínanse
\[
 p_j=\left\lfloor\frac j\delta\right\rfloor,\qquad
 v_j=\mathcal C_{p_j},\qquad h_j=p_{j+1}-p_j.
\]
La desigualdad \(p_j\delta<j<(p_j+1)\delta\) identifica \(p_j\)
como la posición del evento \(j\). Como \(0<\delta<1\), da también
\(\lfloor(p_j+1)\delta\rfloor=j\), y por tanto
\[
 v_j=p_j-j,\qquad
 v_{j+1}-v_j=h_j-1.
\]
Si \(s_j=22-h_j\), entonces \(s_j\in\{0,1\}\) y
\begin{equation}
 [v_{j+1}]_3=[v_j-s_j]_3.
 \label{eq:vacancia-regimen}
\end{equation}
Un intervalo de longitud \(22\) conserva la clase; uno de longitud
\(21\) la desplaza una unidad en sentido negativo. El selector canónico
se aplica a la fase positiva \(r_j=1+[v_j]_9\):
\[
 \tau_j=M_{\rm ph}(r_j).
\]
Queda determinado así un régimen trítico por el índice de capacidad
retenido en cada evento, además de la codificación binaria de la vacancia.

Las posiciones de los intervalos cortos forman la codificación mecánica
inducida de pendiente \(22-1/\delta\). Sus separaciones \(6\) y \(7\)
determinan las longitudes de las mesetas de la clase \([v_j]_3\), una
vez fijados el origen y la convención de extremos. Las primeras clases
son \(2\) durante \(6\) eventos, \(1\) durante \(7\) y \(0\) durante
\(7\); el primer tramo de \(6\) es la restricción inicial de una meseta.
Sus firmas recorren \(0,-1,+1\). Las ecuaciones
\eqref{eq:capacidad-vacancia} y \eqref{eq:vacancia-regimen} explicitan el
enlace entre refinamiento, vacancias y selección del régimen. La lectura
cuadrática del TRIT realiza después esos tipos mediante sus respectivos
generadores; el acontecimiento temporal y su representación geométrica
conservan así una dependencia identificable.

La inversión del índice monótono proporciona otra observación útil:
\[
 \mathcal T(k)=\min\{t:\mathcal C_t=k\}
   =\left\lceil\frac k\rho\right\rceil-1\qquad(k\geq1).
\]
Si \(\sigma=\rho^{-1}-1\), el número de bloques necesarios para ganar
\(m\) unidades de capacidad es
\[
 \mathcal T(k+m)-\mathcal T(k)
 =m+\lceil(k+m)\sigma\rceil-\lceil k\sigma\rceil
 \in\{m+\lfloor m\sigma\rfloor,m+\lceil m\sigma\rceil\}.
\]
En particular, \(9\) unidades de capacidad requieren \(9\) o \(10\)
bloques. El período de una fase finita permanece exacto, mientras el
tiempo de retorno medido por el otro índice admite dos valores. Esta
distinción es parte de la dinámica del refinamiento y no una corrección
de sus valores numéricos a posteriori.
